\subsection{Cascades 3 and 4 fail the rows-to-columns condition}\label{sec:cascade-cuts}

All five cascades are excluded by the regular-recursion obstruction of the second paper~\cite[Theorem~7.5]{SecondPaper2026}, as Astra's
memo~\cite{CascadeObstructionMemo2026} verifies (\cref{prob:cascades}). For every fixed work alphabet there are $\delta>0$ and $d_0$ such that, at
every dyadic $d\ge d_0$, a machine with one of the five traces and states $\{0,\dots,q\}$ that transposes the inputs of a set $S$
satisfies $\clog{(q+1)}+L-\log_2\abs S\ge\delta L$. That argument concerns the motion of the heads and gives no numerical threshold.

For Cascades~3 and~4 a cut does more. Condition~(iv) of the specification already fails for the rows $A_0=\{1,2\}$: an explicit set
of $o(db)$ steps separates their reads from the outputs they feed. Counting single cells instead of whole steps then bounds the work
alphabet and the number of states separately, and both may grow with~$d$. The cut also locates the loss of information.

\begin{definition}[Cell network]\label{def:cellnet}
Let a trace have horizon $T$. For a work tape $j$ and a step $t<T$, let $\mathrm{next}_j(t)$ be the next step after $t$ at which head
$j$ scans the cell it scans at step~$t$, if there is one. The \emph{cell network} of the trace has the vertices $0,\dots,T-1$, a
\emph{state edge} $t\to t+1$ for every $t+1<T$, and a \emph{cell edge} $t\to\mathrm{next}_j(t)$ of tape $j$ for every $j$ and $t$ for which
$\mathrm{next}_j(t)$ exists. A cell edge $t\to t+1$ is a \emph{wait}.
\end{definition}

The cell network runs forward in time and keeps waits as edges of their own. Reversing every edge of the dependency DAG gives edges of
the cell network: a state edge $t\to t-1$ becomes the state edge $t-1\to t$, and the revisit edge of a revisit $(p,u)$ becomes the cell
edge $p\to\mathrm{next}_j(p)=u$.

\begin{lemma}[Cell-level cut-set bound]\label{lem:cellcut}
\notformalized
Let $M$ be oblivious with trace $\tau$, work alphabet $\{0,\dots,g\}$ and states $\{0,\dots,q\}$, and let it transpose at size $d$.
Let $A$ be a set of input positions, $V\subseteq A$, and $C$ a set of edges of the cell network of $\tau$, consisting of $C_{\mathrm{cell}}$
cell edges and $C_{\mathrm{state}}$ state edges, such that every path from a step that reads a position of $A$ to a step of $\Out(V)$
uses an edge of $C$. Then
\[
  \abs V\ \le\ \abs{C_{\mathrm{cell}}}\log_2(g+1)+\abs{C_{\mathrm{state}}}\log_2(q+1).
\]
\end{lemma}

\begin{informalproof}
Fix the input bits outside $A$. Let $U$ be the set of steps reachable by a path that avoids $C$ from a step that reads a position of
$A$. Every edge from a step of $U$ to a step outside $U$ lies in $C$. A step reads the state left by the step before it and the
symbols last written on the cells it scans (the blank on a first visit), so the state it enters, the symbols it writes and the bit it
appends are a function of the values carried by its incoming edges and of the input bit it reads. By induction on $t$, everything a
step $t\notin U$ computes is a function of the values carried by the edges of $C$: such a step reads no position of $A$, and each of
its incoming edges starts outside $U$ or lies in $C$. The steps of $\Out(V)$ lie outside $U$, and at $\Out(v)$ the machine appends
the input bit $v$. So the bits $(x_v)_{v\in V}$ are a function of the values on $C$. They take all $2^{\abs V}$ values as the bits of
$A$ vary, and the values on $C$ take at most $(g+1)^{\abs{C_{\mathrm{cell}}}}(q+1)^{\abs{C_{\mathrm{state}}}}$.
\end{informalproof}

\begin{theorem}[Cascades 3 and 4]\label{thm:cascade-cuts}
\notformalized
There are absolute constants $c$ and $d_0$ with the following property. Let $d=2^k\ge d_0$, let $\tau$ be the trace of Cascade~3 or
of Cascade~4 at size $d$, and put $A_0=\{1,2\}$, $A_1=\{0,\dots,d-1\}$, and $\nu=(db)^{1/5}$ for Cascade~3 and $\nu=(db)^{1/6}$ for
Cascade~4.
\begin{enumerate}[label=(\alph*)]
\item Some set of at most $c\,db/\nu$ steps separates $\Out(\Qstar(A_0,A_1))$ from $\Lin(A_0)$. Hence, by \cref{lem:R}, no oblivious
machine with trace $\tau$ and $\wval+1<\nu/c$ transposes at size~$d$.
\item No oblivious machine with trace $\tau$, work alphabet $\{0,\dots,g\}$ and states $\{0,\dots,q\}$ with
$\log_2(g+1)+\log_2(q+1)<\nu/c$ transposes at size~$d$.
\end{enumerate}
\end{theorem}

\begin{proof}[Proof outline]
The proof is in \cref{app:cascade-cuts}. Both schedules read rows~$1$ and~$2$ only while loading the input, at the $2db$ load steps
that read the positions $db,\dots,3db-1$. Call a step tainted if it is reachable from these reads in the cell network, with three
exceptions: taint never passes along a state edge; it passes along a wait only within one step (two steps on Cascade~4) of the head's
arrival on the cell; and a set $C_0$ of cell edges is cut, namely the tape-$1$ edges that leave tainted steps of the transfer at one level
$\rho$ and next visit their cell after it. Every edge from a tainted step to an untainted step then lies in $C_0$, or
leaves the last step of a maximal interval of tainted steps, a \emph{run} (\cref{lem:closure}). So the cut costs $\abs{C_0}$ plus a
bounded number of edges per run, and the tainted steps of $\Out(\Qstar)$ are the outputs given up.

Each fold of these schedules starts with a transfer of half of its source, in which tape~$1$ moves at half speed while tape~$0$ crosses
the first two thirds of that half. Rows~$1$ and~$2$ stay in that part at every level, so their window halves at each level: at level
$\rho$ it has at most $db/2^{\rho-1}+O(1)$ cells, and the transfer squeezes it into $db/2^\rho+O(1)$ cells of tape~$1$. Cutting the edges
out of these cells leaves every deeper level untainted. Below $\rho$, the backward stage
spreads the window again, by a bounded factor per level. At the level we choose it has $O(6^\rho)$ runs on Cascade~3 and $O(11^\rho)$
on Cascade~4, and the final sweep gives up $O(6^\rho b)$, respectively $O(11^\rho b)$, outputs. It also gives up the $4b$ outputs of
the entries $(1,1)$, $(1,2)$, $(2,1)$ and $(2,2)$, at whose append steps it reads the loaded copy of rows~$1$ and~$2$. Choosing $2^\rho$
near $\nu$ makes the cut and the outputs given up both $o(db)$. For~(a), the tails of the cut edges together with the roots given up
separate in the dependency DAG; for~(b), \cref{lem:cellcut} applies to the at least $db$ outputs that the cut separates.
\end{proof}

\begin{remark}[Other entry widths]\label{rem:cascade-widths}
The geometric part of the proof holds for every $b\ge1$ once $d\ge32$. Choosing $2^\rho$ near $d^{1/5}$ (Cascade~3) or $d^{1/6}$
(Cascade~4) instead, the count of \cref{lem:cellcut} excludes the two traces, for $d\times d$ matrices with $b$-bit entries and any
$b\ge1$, once $\log_2(g+1)+\log_2(q+1)$ is below $d^{1/5}/c$, respectively $d^{1/6}/c$. In particular it excludes them for matrices of
bits.
\end{remark}

\begin{remark}[Rows~0 and~1]\label{rem:cascade-rows}
\Cref{lem:R} counts every step at which the input head is in rows $A_0$. Both schedules keep the input head on bit~$0$ throughout their
initial positioning, a linear number of steps whose writes later steps revisit, so with $A_0=\{0,1\}$ these steps join $\Lin(A_0)$.
The choice of rows in \cref{lem:R} is free, and rows~$1$ and~$2$ avoid them.
\end{remark}

The two exclusions of Cascades~3 and~4 differ in scope. The memo~\cite{CascadeObstructionMemo2026} covers all five cascades, and it excludes every
machine with $\clog{(q+1)}<\delta L$ that is correct on more than a $2^{\clog{(q+1)}-\delta L}$ fraction of the inputs; but it needs a
fixed work alphabet, and its $\delta$ and $d_0$ depend on the alphabet and are obtained by contradiction. \Cref{thm:cascade-cuts}
covers two cascades and machines that transpose every input, but it allows alphabets and controllers of up to about $\nu$ bits, with
absolute constants.

It also names the mechanism that fails. The half-speed transfer squeezes the data of rows~$1$ and~$2$ into half as many cells at
each level. On Cascade~3, after about $\frac15\log_2(db)$ levels these data pass through about $(db)^{4/5}$ cells, where the backward
stage is still cheap to cut. Cascades~1 and~2 use the same transfer, but before their cascade a shuffle, in which both work heads read
and write at every step, spreads each row over several distant intervals of the tape, so that the data of a pair of rows is no longer
confined to the squeezed part. We
know no cut that excludes Cascades~1, 2 or~5.
